\documentclass[11pt]{report}
\usepackage[a4paper,margin=1in]{geometry}
\usepackage{amsmath,amssymb}
\usepackage{graphicx}
\usepackage{booktabs}
\usepackage{longtable}
\usepackage{array}
\usepackage[colorlinks=true,linkcolor=blue,urlcolor=blue]{hyperref}
\usepackage{fancyhdr}
\usepackage{xcolor}
\usepackage{titlesec}
\usepackage{enumitem}
\usepackage{calc}

\definecolor{hbblue}{HTML}{2A4D8F}
\definecolor{warnred}{HTML}{9C2B2B}
\definecolor{boxgrey}{HTML}{F4F4F6}

\titleformat{\chapter}[display]{\bfseries\Large\color{hbblue}}{\chaptertitlename\ \thechapter}{10pt}{\huge}
\pagestyle{fancy}
\fancyhf{}
\fancyhead[R]{\small\itshape Pinhole Sunglasses --- Physics and Fabrication Guide}
\fancyfoot[C]{\thepage}

\newsavebox{\kbsave}
\newsavebox{\wbsave}

\newenvironment{keybox}[1]%
{\par\vspace{8pt}\begin{lrbox}{\kbsave}\begin{minipage}{0.88\textwidth}%
{\bfseries\color{hbblue}#1}\par\vspace{4pt}\ignorespaces}%
{\end{minipage}\end{lrbox}%
\begin{center}\setlength{\fboxrule}{1.2pt}\setlength{\fboxsep}{8pt}%
\fcolorbox{hbblue}{boxgrey}{\usebox{\kbsave}}\end{center}\vspace{8pt}\par}

\newenvironment{warnbox}[1]%
{\par\vspace{8pt}\begin{lrbox}{\wbsave}\begin{minipage}{0.88\textwidth}%
{\bfseries\color{warnred}#1}\par\vspace{4pt}\ignorespaces}%
{\end{minipage}\end{lrbox}%
\begin{center}\setlength{\fboxrule}{1.2pt}\setlength{\fboxsep}{8pt}%
\fcolorbox{warnred}{white}{\usebox{\wbsave}}\end{center}\vspace{8pt}\par}

\newcommand{\ahole}{a_{\text{hole}}}
\newcommand{\dgap}{d_{\text{gap}}}
\newcommand{\vr}{v_r}
\newcommand{\vf}{v_f}
\newcommand{\Dpupil}{D_{\text{pupil}}}

\title{\textbf{\Huge Pinhole Sunglasses}\\[0.3cm]
\textcolor{hbblue}{\huge Physics and Fabrication Guide}\\[0.5cm]
\large\itshape From wave optics to a brass shim you can etch:\\
every equation, why it is there, and what to do at the bench}
\author{}
\date{IYPT 2027 Problem 1 --- working reference}

\begin{document}
\maketitle
\tableofcontents

\chapter*{How to read this document}
\addcontentsline{toc}{chapter}{How to read this document}

Part~I builds the physics from the ground up: why a pinhole helps a myopic eye at all,
why making it smaller eventually stops helping, and how those two competing effects are
computed exactly rather than estimated. Part~II turns the resulting design numbers into
a mask you can print and a shim you can etch, with the specific checks that catch the
failure modes of photochemical etching before you waste material.

Every number quoted here is read from this project's own output files, not invented for
illustration. The two source files that matter are:

\begin{itemize}[leftmargin=1.2cm]
\item \texttt{array\_data/evaluation\_results/gaze\_weighted\_winner\_per\_M.csv} --- the final designs
\item \texttt{fabrication\_package/fabrication\_summary.csv} --- the mask dimensions derived from them
\end{itemize}

These two files agree row-for-row. That is worth knowing before you start: the thing you
are about to fabricate \emph{is} the thing the model selected, not an approximation to it.

\part{The Physics}

\chapter{Why a pinhole helps, and why smaller is not better}

\section{The core tension in one paragraph}

A myopic eye focuses light in front of the retina. Every point in the scene therefore lands
on the retina as a disc rather than a point, and the disc's size is set by how wide a cone
of light the eye accepts. Narrow that cone with a pinhole and the disc shrinks --- this is
pure geometry and it is why squinting works. But light is a wave, and forcing it through a
small hole makes it spread by diffraction. Below some aperture the diffraction spreading
grows faster than the geometric blur shrinks, and the image gets worse again.

\begin{keybox}{The governing trade-off}
Geometric (defocus) blur $\propto$ aperture size. \\
Diffraction blur $\propto$ 1 / aperture size. \\
Somewhere between them is a minimum. Finding it, correctly, is the entire single-hole problem.
\end{keybox}

\section{Two ways to compute it, one of which is wrong}

The tempting approach is to write down a formula for each blur and combine them:
\begin{equation}
r_{\text{geo}} = \frac{a_{\text{stop}}}{2}\cdot\frac{|\vr-\vf|}{\vf},
\qquad
r_{\text{diff}} = \frac{1.22\,\lambda\,\vr}{a_{\text{stop}}},
\qquad
r_{\text{total}} \overset{?}{=} \sqrt{r_{\text{geo}}^2+r_{\text{diff}}^2}
\end{equation}

That last step --- root-sum-square (RSS) --- is physically wrong. Two blur mechanisms combine
by \emph{convolving their point-spread functions}, not by adding their radii in quadrature.
RSS is a heuristic borrowed from independent Gaussian error terms, and neither blur here is
Gaussian: defocus produces a roughly uniform disc, diffraction produces an Airy pattern with
rings extending well beyond the central lobe.

\begin{warnbox}{How wrong, measured}
This project quantified the RSS error against the exact numerical result rather than asserting
it: \textbf{64.8\% mean divergence, up to 192.8\% maximum}, at the retina plane, worst at mild
myopia. That is a factor of two to three, not a rounding difference. Do not defend RSS; report
that you measured its failure.
\end{warnbox}

The correct approach does not separate the two blurs at all. It propagates the actual light
field through the actual geometry, and both effects emerge together because they were never
distinct physical processes to begin with --- they are two descriptions of the same wave.

\chapter{Modelling the myopic eye}

\section{Axial myopia: the retina moves, the lens does not}

Myopia can arise two ways: the eye's optical power is too strong (refractive myopia), or the
eyeball is too long (axial myopia). Real short-sightedness is overwhelmingly the second. This
project models it that way:
\begin{equation}
L_0 = \frac{n'}{P_0}, \qquad L_{\text{ax}} = \frac{n'}{P_0-M}
\label{eq:axial}
\end{equation}
with $P_0 = 60$\,D the emmetropic eye power (held \emph{fixed} at every myopia level),
$n' = 1.336$ the vitreous refractive index, and $M$ the myopia in dioptres.

\textbf{Worked example, $M=4$\,D:}
\begin{align}
L_0 &= 1.336/60 = 22.267\ \text{mm} \quad\text{(where the light focuses)}\\
L_{\text{ax}} &= 1.336/56 = 23.857\ \text{mm} \quad\text{(where the retina actually is)}\\
\Delta L &= 1.59\ \text{mm}
\end{align}

So a $-4$\,D eye is a normal eye with its retina sitting 1.59\,mm too far back. The focal
length is unchanged. That $\Delta L$ is the entire disease, and it is the reason the next
chapter's propagation distance matters so much.

\begin{keybox}{Why this convention was chosen}
Holding $P_0$ fixed and growing the axial length is the standard clinical model of \emph{simple
axial myopia}. It also means one number ($M$) fully specifies the geometry, so every result can
be reported per-dioptre without ambiguity about which parameter was varied.
\end{keybox}

\section{The single most important architectural point}

Once the light passes the eye's lens, you must propagate it to the retina. The obvious move ---
multiply by the lens's quadratic phase, then take a Fourier transform --- is a trap.

A single FFT after a lens phase gives you the field at \textbf{the lens's own back focal plane}.
For a healthy eye that is the retina, and the answer is right. For a myopic eye the retina is
\emph{not} at the focal plane --- that mismatch is the definition of myopia --- and the FFT
silently computes the field at the wrong location.

\begin{warnbox}{What the bug produces}
A beautiful, sharp Airy pattern computed at the eye's true focal point: the one plane in the
whole system where defocus blur is exactly zero. You would be modelling a perfectly corrected
eye and calling it myopic. The plot looks plausible; the physics is inverted.
\end{warnbox}

\textbf{The fix:} replace the FFT with a second Angular Spectrum propagation over the actual
physical lens-to-retina distance $\vr$. This is a one-line change with three consequences:

\begin{enumerate}[leftmargin=1.2cm]
\item Geometric defocus blur is no longer computed separately --- it appears automatically,
because you propagated a wavefront curved for the wrong focal length across the correct distance.
\item The RSS combination is not merely inaccurate, it is \emph{structurally unnecessary}: the two
blurs were never separated, so there is nothing to recombine.
\item The earlier ``convolve an Airy PSF with a defocus disc'' step is superseded --- that was an
approximation to exactly what this propagation now computes directly.
\end{enumerate}

\chapter{The wave-optics engine}

\section{The Angular Spectrum Method, explained physically}

Any field at a plane can be decomposed into plane waves travelling in different directions.
Each plane wave, crossing a distance $z$, simply acquires a phase. Reassemble them and you have
the field at the new plane. That is the whole idea:
\begin{equation}
U(x,y,z) = \mathcal{F}^{-1}\Big\{\mathcal{F}\{U_0(x,y)\}\cdot H(f_x,f_y;z)\Big\},
\qquad
H = \exp\!\Big(ikz\sqrt{1-\lambda^2(f_x^2+f_y^2)}\Big)
\label{eq:asm}
\end{equation}
valid for propagating components $f_x^2+f_y^2 \le 1/\lambda^2$; beyond that the square root
turns imaginary and those components decay exponentially (evanescent waves), so $H=0$.

This is the \textbf{exact} solution of the scalar Helmholtz equation between parallel planes.
It makes no near-field or far-field assumption, which is precisely why it is the right default
engine for this problem: the pinhole-to-cornea gap and the lens-to-retina distance sit in
different diffraction regimes, and switching between Fresnel and Fraunhofer closed forms at
each stage would introduce approximation errors that ASM simply does not have.

\section{Which regime you are in: the Fresnel number}
\begin{equation}
N_F = \frac{R_{\text{hole}}^2}{\lambda z}, \qquad R_{\text{hole}} = \ahole/2
\label{eq:NF}
\end{equation}
$N_F \ll 1$: far field (Fraunhofer), a clean Airy pattern. $N_F \sim 1$--$10$: Fresnel transition,
several comparably bright rings. $N_F \gg 10$: nearly geometric, the hole projects a sharp-edged
patch.

\textbf{Worked example.} $\ahole = 300\,\mu$m so $R = 150\,\mu$m, $\lambda = 550$\,nm,
$\dgap = 12$\,mm:
\begin{equation}
N_F = \frac{(1.5\times10^{-4})^2}{(5.5\times10^{-7})(1.2\times10^{-2})}
= \frac{2.25\times10^{-8}}{6.6\times10^{-9}} \approx 3.41
\end{equation}
Squarely in the transition regime. Expect multiple rings, not a textbook Airy disc.

Because this project's sweeps span $N_F$ from well under 1 to over 20, no single closed form
covers the range --- another reason ASM is used throughout rather than as a fallback.

\section{Band-limiting: a numerical necessity, not physics}

A finite simulation window with finite grid spacing cannot represent arbitrarily steep
propagation angles. A plane wave steep enough to exit the window during propagation will,
in a discrete FFT, wrap around and re-enter from the opposite side --- appearing as spurious
fine ring structure that looks exactly like real diffraction.

The Matsushima band-limit suppresses spectral content beyond the angle at which this happens.
This is anti-aliasing, applied deliberately and with the suppressed energy fraction logged
per run. It is not a physical loss mechanism, and you should say so if asked.

\begin{keybox}{Why the energy-conservation check had to be rewritten}
Since $H$ has unit modulus on every kept component and is exactly zero elsewhere, the apparent
energy loss is \emph{identically} the band-limit-suppressed fraction:
$\varepsilon_{\text{energy}} \equiv f_{\text{suppressed}}$ --- confirmed against real sweep data at
correlation $1.000000$, standard deviation $6\times10^{-14}$. Testing $\varepsilon$ against a
$0.1\%$ tolerance was therefore guaranteed to fail on every row, not from any defect but from
re-flagging intentional, already-tolerated suppression. The fix tests the \emph{residual}
$\varepsilon_{\text{energy}} - f_{\text{suppressed}}$, restoring the check's ability to detect a
genuine bug.
\end{keybox}

\chapter{The full pipeline, stage by stage}

\begin{center}
\fbox{\parbox{0.92\textwidth}{\centering
plane wave $\rightarrow$ pinhole mask $U_0$ $\rightarrow$ \textbf{ASM}$(\dgap)$ $\rightarrow$
pupil clip $\rightarrow$ lens phase $\rightarrow$ \textbf{ASM}$(\vr)$ $\rightarrow$
$I_{\text{retina}} = |U|^2$}}
\end{center}

\begin{description}[leftmargin=0.8cm,style=nextline]
\item[1. Source and mask.] A distant object gives a collimated (plane) wave. The pinhole is a
binary amplitude mask: 1 inside the hole, 0 outside.

\item[2. ASM across the gap $\dgap$.] Typically 12--15\,mm, the spectacle vertex distance. This is
where the light diffracts \emph{before} it ever reaches the eye.

\item[3. Pupil clipping.] Multiply by a disc of diameter $\Dpupil$ (default 4000\,$\mu$m). For gaze
studies this disc is laterally offset --- see Chapter~6.

\item[4. Lens phase.] Multiply by
$\exp\!\big(-i\frac{k}{2f}(x^2+y^2)\big)$
using the eye's \emph{actual} focal length. The eye is modelled as it is; the pinhole is the
only correction present.

\item[5. ASM across $\vr$.] The corrected stage. $\vr = L_{\text{ax}}$ from Eq.~\ref{eq:axial}.

\item[6. Intensity.] $I(x,y) = |U(x,y,\vr)|^2$, then radial integration to get encircled energy.
\end{description}

\begin{warnbox}{A modelling subtlety that is easy to get backwards}
The aperture that appears in any diffraction formula at the eye is \textbf{the propagated light
spot at the cornea plane}, not the physical hole diameter. The hole diffracts across $\dgap$
before the pupil sees it; by then the illuminated patch is generally wider than the hole. Using
the physical hole diameter in a cornea-plane formula understates the spot and will bias any
hand-check you run against the pipeline.
\end{warnbox}

\section{Validating the engine before trusting any output}

Set $\dgap \to 0$ and $M = 0$. The whole chain must then reduce to the textbook Fraunhofer result:
\begin{equation}
I(r) = I_0\left(\frac{2J_1(x)}{x}\right)^2,\quad x = \frac{\pi a r}{\lambda \vr},
\qquad
EE(x) = 1 - J_0^2(x) - J_1^2(x)
\label{eq:airy}
\end{equation}

It does, to 2\% on $r_{EE}(50\%)$ and $r_{EE}(90\%)$.

$r_{EE}(80\%)$ needs an 8\% tolerance, and the reason is worth understanding rather than excusing:
at the first null $x = 3.8317$, $EE \approx 0.838$. The 80\% radius therefore sits almost exactly
on the first dark ring, where the encircled-energy curve is momentarily flat. On a flat curve, a
tiny error in energy maps to a large error in radius. Confirmed empirically: grid sizes from
1024 to 8192 --- an order of magnitude more compute --- never brought this percentile below
2.3--4.6\% error. It is ill-conditioning, not a bug, and no amount of resolution fixes it.

\begin{keybox}{Also note what Eq.~\ref{eq:airy} tells you about blur metrics}
At the first null only 83.8\% of the energy is enclosed. Over 16\% of the light sits outside the
``obvious'' central disc. Any description of blur by a null radius alone therefore misrepresents
where the light actually is --- which is why this project reports $r_{EE}(P)$ at
$P = 0.5, 0.7, 0.8, 0.9$ rather than a single number.
\end{keybox}

\chapter{Blur metrics and where they break}

You will be asked which number measures ``blur.'' There is no single safe answer, and knowing
the failure modes is more valuable than picking one.

\begin{longtable}{p{3.0cm}p{5.5cm}p{5.5cm}}
\caption{Blur metrics: what each measures and how it fails.}\\
\toprule
\textbf{Metric} & \textbf{Measures} & \textbf{Failure mode} \\
\midrule
\endhead
$r_{EE}(50\%)$ & Radius enclosing half the energy & Stable and monotonic, but overestimates
clinical acuity loss by roughly $7\times$ against a real clinical study \\
FWHM & Width of the central core only & Tracks the same clinical study to $\sim1.6\times$, but
\textbf{ring-hops}: jumped $4.7 \to 154.7\,\mu$m between apertures 250\,$\mu$m apart at $M=4$\,D \\
Strehl ratio & Peak intensity vs.\ diffraction limit & Genuinely multi-peaked: 6--7 local maxima,
swings up to $8.4\times$; finer sampling revealed \emph{more} structure, not less \\
MTF50 & Spatial frequency at 50\% contrast & Most informative, but is a curve, not a scalar \\
\bottomrule
\end{longtable}

\begin{keybox}{How to state this without contradicting yourself}
Two sentences in this project's own record appear to conflict: ``FWHM is the more clinically
accurate metric'' and ``FWHM is unreliable at the retina plane.'' Both are true at different
scales. The resolution to state out loud:

\emph{Per-aperture, FWHM is unusable --- it discontinuously jumps between diffraction rings.
Across a coarse sweep its trend is closer to clinical reality than $r_{EE}(50\%)$, which is
smooth and stable but biased pessimistic by about $7\times$. We therefore report $r_{EE}(50\%)$
as the optimisation metric and state its bias explicitly, rather than optimising against a
metric that is discontinuous in the design variable.}
\end{keybox}

\section{Why Strehl's oscillation is real physics, not numerical noise}

On-axis intensity behind a circular aperture oscillates with Fresnel number:
\begin{equation}
I(0,z) = 4I_0\sin^2\!\left(\frac{\pi N_F}{2}\right)
\label{eq:onaxis}
\end{equation}
swinging between 0 and $4I_0$ as successive Fresnel zones alternately reinforce and cancel on
axis. This project's simulated data matches that closed form at correlation \textbf{0.997}. So
when Strehl ratio oscillates non-monotonically with aperture, that is Fresnel-zone interference
doing exactly what Eq.~\ref{eq:onaxis} says it must --- and it is also why the deprecated
quadrature estimate $a_{\text{eff}} \approx \sqrt{\ahole^2 + (2.44\lambda\dgap/\ahole)^2}$, which
is smooth and monotonic by construction, cannot represent this regime.

\chapter{The three apertures, reconciled}

Three different aperture numbers appear across this project. They are not three competing
answers. They are the same blur-versus-aperture curve read under three different constraint sets.

\begin{longtable}{p{4.2cm}p{3.0cm}p{6.5cm}}
\caption{The same optimisation, three constraint sets.}\\
\toprule
\textbf{Question asked} & \textbf{Answer} & \textbf{What is binding} \\
\midrule
\endhead
Minimise blur only & 300--400\,$\mu$m & Diffraction vs.\ defocus alone \\
Minimise blur subject to a retinal brightness floor & 900--1250\,$\mu$m & Illuminance constraint;
lands inside the clinical bracket 0.94--1.75\,mm \\
Minimise blur subject to brightness \emph{and} array ghost fusion & 530\,$\mu$m
(Shen's universal design) & Ghost-fusion criterion in a dense-array regime \\
\bottomrule
\end{longtable}

\section{The brightness constraint, from energy conservation}
\begin{equation}
E_{\text{rel}} = \left(\frac{\min(\ahole,\Dpupil)}{\Dpupil}\right)^2
\label{eq:erel}
\end{equation}
Note what this does \emph{not} use: any diffraction-broadened spot size. Diffraction redistributes
power spatially; it does not create or destroy it. Retinal illuminance must therefore be computed
from the raw physical aperture area. This is a first-principles result, not a model choice.

\textbf{Worked example.} $\ahole = 1000\,\mu$m against a 4\,mm pupil:
$E_{\text{rel}} = (1000/4000)^2 = 0.0625 = 6.25\%$. This is why $T_{\min}=6.25\%$ appears as a
sweep value --- it is exactly the 1\,mm aperture's transmission, not an arbitrary round number.

\begin{longtable}{ccc}
\caption{Constrained optimum versus brightness floor.}\\
\toprule
$T_{\min}$ & Constrained optimum & vs.\ clinical bracket \\
\midrule
\endhead
0\% & 300--400\,$\mu$m & below \\
1\% & 400\,$\mu$m & below \\
3\% & 700\,$\mu$m & below \\
6.25\% & 1000\,$\mu$m & \textbf{inside} \\
9\% & 1250\,$\mu$m & \textbf{inside} \\
12.25\% & 1500\,$\mu$m & upper edge \\
\bottomrule
\end{longtable}

\begin{keybox}{The finding, stated in one sentence}
A brightness floor of 6--9\% of full natural-pupil illuminance --- physically modest, comparable
to ordinary sunglass tinting --- moves the predicted optimum from 300--400\,$\mu$m into the
verified clinical range of 0.94--1.75\,mm, uniformly across every myopia level tested. The
diffraction physics was never wrong; blur-only optimisation was answering the wrong question.
\end{keybox}

\chapter{From one hole to an array}

Three quantities govern array design. Keep them separate --- they fail independently, and
collapsing them into one score hides which one broke.

\section{Viability: does light get through at all}

The pupil and a hole are both circles. Light passes if they overlap:
\begin{equation}
d - R_h < R_p
\label{eq:viability}
\end{equation}
where $d$ is the hole-centre-to-pupil-centre distance, $R_h$ the hole radius, $R_p$ the pupil
radius. If they do overlap, the illuminated area is the circle--circle intersection:
\begin{equation}
\begin{split}
A = R_p^2\cos^{-1}\!\left(\frac{d^2+R_p^2-R_h^2}{2dR_p}\right)
  &+ R_h^2\cos^{-1}\!\left(\frac{d^2+R_h^2-R_p^2}{2dR_h}\right)\\
  &- \tfrac12\sqrt{(-d+R_p+R_h)(d+R_p-R_h)(d+R_h-R_p)(d+R_p+R_h)}
\end{split}
\label{eq:overlap}
\end{equation}

\textbf{At straight-ahead gaze} ($\varphi=0$) the pupil is centred on the optical axis, so
$d = \sqrt{x^2+y^2}$ --- just the hole's distance from the origin.

\textbf{At non-zero gaze} the eye rotates about a centre roughly $R_{\text{rot}} \approx 13$--14\,mm
behind the pupil, so the pupil centre itself translates:
\begin{equation}
\Delta(\varphi) = R_{\text{rot}}\sin\varphi
\label{eq:gazeshift}
\end{equation}
and $d$ becomes the hole-to-\emph{shifted}-pupil-centre distance. This is why a design can be
perfectly viable looking straight ahead and go dark at $15^\circ$ --- and why the final model in
this project is gaze-weighted rather than on-axis.

\section{Ghosting: two lit holes, two images}

If two holes are simultaneously inside the pupil, each sends its own beam through the lens at a
different lateral position. In a defocused eye those beams land at different retinal points:
\begin{equation}
\Delta r = \frac{s\,M}{P_0-M}
\label{eq:ghost}
\end{equation}
with $s$ the centre-to-centre hole separation. This project validated Eq.~\ref{eq:ghost} against
the full model at $R^2 = 0.995$.

\begin{keybox}{The sanity check that caught a $\sim100\times$ formula error}
$\Delta r \to 0$ as $M \to 0$. An in-focus eye converges every beam to the same point regardless
of where it entered the pupil --- so a healthy eye cannot ghost, by construction. Any candidate
ghost formula that does not vanish at $M=0$ is wrong. An earlier version of this formula failed
exactly that test.
\end{keybox}

A pair is tolerable if $\Delta r$ falls below the eye's fusion threshold $c_{\min}$ (or the
measured total blur $c_{\text{tot}}$, whichever is larger) --- the two images blur into one.

\section{Packing density}

Geometry, not physics, and the place an earlier bug lived. Family-specific unit-cell formulas
disagreed with hull-based ones by 1.5--2.25$\times$. The family-agnostic version:
\begin{equation}
A_{\text{array}} = A_{\text{hull}} + P_{\text{hull}}\cdot R_h + \pi R_h^2
\label{eq:density}
\end{equation}
the convex hull of hole centres, buffered outward by one hole radius --- because a hole's
physical extent sticks out past its own centre point.

\chapter{The design this project actually selected}

\section{Reading the final table}

From \texttt{gaze\_weighted\_winner\_per\_M.csv}:

\begin{longtable}{cccccc}
\caption{Gaze-weighted winners. Pitch is 4000\,$\mu$m and rotation $0^\circ$ in every row.}\\
\toprule
$M$ (D) & Family & $N$ & $\ahole$ ($\mu$m) & Lit holes at $\varphi{=}0$ & Time-wtd.\ ghost \% \\
\midrule
\endhead
0 & square & 20 & 2000 & 1 & 0.0 \\
1 & hexagonal & 20 & 1250 & 1 & 0.0 \\
2 & hexagonal & 20 & 750 & 1 & 0.0 \\
3 & hexagonal & 20 & 750 & 1 & 0.0 \\
4 & hexagonal & 20 & 600 & 1 & 0.0 \\
5 & hexagonal & 20 & 600 & 1 & 0.0 \\
6 & hexagonal & 100 & 500 & 1 & 0.0 \\
8 & hexagonal & 50 & 400 & 1 & 0.0 \\
global & square & 20 & 2000 & 1 & 0.0 \\
\bottomrule
\end{longtable}

Two columns matter more than the rest. \textbf{Lit holes at $\varphi=0$ is 1 in every row}, and
\textbf{worst-case ghost margin is \texttt{inf} in every row} --- infinite because there are no
simultaneously-lit pairs to evaluate.

\begin{keybox}{The design principle, named}
Pitch (4000\,$\mu$m) equals the pupil diameter (4000\,$\mu$m). The array is deliberately sparse
enough that \textbf{only one hole illuminates the pupil at a time}. Each hole acts as an
independent single pinhole; the array exists to provide field of view as the eye rotates, not to
put multiple apertures in front of the pupil at once.

This is a fundamentally different architecture from a dense array (e.g.\ Shen's $s = 1.10$\,mm
against a 2--4\,mm pupil, which puts several holes in the pupil simultaneously and must then rely
on ghost \emph{fusion} to stay usable).
\end{keybox}

\section{Why \texorpdfstring{$c_{\min}$}{c\_min} does not bind your result}

$c_{\min} = 55\,\mu$m is inherited from Shen and is that paper's weakest parameter --- its own
authors describe it as a single-parameter qualitative calibration anchored to one clinical
geometry, reproducing threshold behaviour only qualitatively (predicted ghost onset $\approx1.1$\,D
against a clinical $\approx4$\,D), and they name direct two-pinhole calibration as the most
valuable falsifiable test of their model. You do not have the resources for that experiment.
You do not need it, and here is the arithmetic proving why.

Take the $M=8$\,D design: hexagonal, 50 holes, $s = 4000\,\mu$m (verified by direct
nearest-neighbour computation over the actual generated coordinates). If two of those holes were
ever lit simultaneously, Eq.~\ref{eq:ghost} gives
\begin{equation}
\Delta r = \frac{4000 \times 8}{60-8} = 615.4\,\mu\text{m}
\end{equation}
against $c_{\min} = 55\,\mu$m. That is $11\times$ over threshold --- catastrophic ghosting, and it
would be catastrophic at $c_{\min} = 210\,\mu$m too, or at any plausible value.

\begin{keybox}{The defensible statement}
This design does not survive ghosting by \emph{fusing} ghost images. It survives by \textbf{never
producing a pair}: with pitch equal to pupil diameter, only one aperture illuminates the pupil at
any gaze direction, so no ghost pair exists to evaluate. The ghost margin is reported as
\texttt{inf} for exactly this reason.

$c_{\min}$ is therefore \textbf{non-binding} on the selected design. Its uncertainty is real and
inherited, but it can only affect designs in the dense-array regime ($s < \Dpupil$), which this
one is not. That is a structural argument from geometry, not a calibration claim, and it does not
require the two-pinhole experiment to make.
\end{keybox}

\begin{warnbox}{The one test that would make this airtight, and costs no lab time}
Re-run the gaze-weighted evaluator at $c_{\min} = 55, 110, 210, 400\,\mu$m with everything else
held fixed (lattice set, aperture range, pupil diameter, gaze weighting). If the winners are
unchanged --- which the argument above predicts --- you can state \emph{``we tested how much this
parameter matters and it does not change our result,''} which is strictly stronger than
\emph{``we used the published value.''} Half a day of compute. Do it before the presentation.
\end{warnbox}

\section{Honest limits of this architecture}

State these yourself rather than let them be extracted:

\begin{itemize}[leftmargin=1.2cm]
\item \textbf{Blind fraction.} Only the $M=6$ and $M=8$ designs show non-zero time-weighted blind
fraction (0.12\% and 0.33\%). Small, but non-zero: at some gaze angles no hole is lit.
\item \textbf{Light throughput is low by construction.} One hole per pupil means transmission is
that hole's area over the pupil area --- for $M=8$ ($400\,\mu$m), $(400/4000)^2 = 1\%$. This is
dark. It is a genuine trade-off against the dense-array approach, and you should present it as
one rather than as a free win.
\item \textbf{No measured acuity change yet.} The chain from ``array reduces simulated blur'' to
``array measurably improves acuity in dioptres'' is not yet closed experimentally.
\end{itemize}

\part{Fabrication}

\chapter{From design number to mask dimension}

\section{The chain, and where each number comes from}

\begin{center}
\fbox{\parbox{0.92\textwidth}{\centering
gaze-weighted winner ($\ahole$, pitch, $N$) $\rightarrow$ hole coordinates $\rightarrow$
undercut compensation $\rightarrow$ mask aperture $\rightarrow$ printed transparency $\rightarrow$
photoresist $\rightarrow$ etch $\rightarrow$ measured hole}}
\end{center}

\texttt{fabrication\_summary.csv} traces row-for-row to \texttt{gaze\_weighted\_winner\_per\_M.csv}.
The fabricated article \emph{is} the selected design --- confirm this yourself before the
presentation, because it is the question you will be asked.

\section{Undercut: why the mask hole is smaller than the hole you want}

Wet chemical etching is isotropic. Ferric chloride attacks the brass in every direction at the
same rate, so it eats sideways under the resist edge as fast as it eats downward. The hole that
emerges is therefore \emph{wider} than the resist opening that defined it.

For a double-sided (simultaneous two-face) etch through a shim of thickness $t$, each face etches
roughly halfway through, and the lateral undercut on each face scales with that depth. The design
rule used here:
\begin{equation}
d_{\text{mask}} = d_{\text{target}} - C_{\text{undercut}},
\qquad
C_{\text{undercut}} \in \{0.5t,\ 1.0t,\ 1.5t\}
\label{eq:undercut}
\end{equation}
with $t = 50\,\mu$m (0.05\,mm brass shim), giving a bracket of 25 / 50 / 75\,$\mu$m compensation.

\begin{warnbox}{This is an uncalibrated design rule, not a measurement}
There is \textbf{no empirical etch-rate or undercut calibration} for this specific brass alloy,
thickness, etchant concentration, temperature, or agitation. The $0.5$--$1.5t$ bracket is a
standard published starting rule. The bracket exists precisely because the true value is unknown
--- which is why the test coupon in Chapter~11 is not optional.
\end{warnbox}

\textbf{Worked example, $M = 8$\,D design.} Target etched diameter $400\,\mu$m:
\begin{align}
d_{\text{mask,nominal}} &= 400 - 1.0(50) = 350\,\mu\text{m} \quad\text{(the active mask set)}\\
d_{\text{mask,low}} &= 400 - 0.5(50) = 375\,\mu\text{m} \quad\text{(least aggressive)}\\
d_{\text{mask,high}} &= 400 - 1.5(50) = 325\,\mu\text{m} \quad\text{(most aggressive)}
\end{align}

Note the naming: ``low'' and ``high'' refer to the assumed \emph{undercut}, so the low-undercut
mask has the \emph{larger} aperture. Getting this backwards will cost you a shim.

\section{The mask files, and what the four variants mean}

Each design generates four SVGs:
\begin{center}
\texttt{M8.0\_positive\_top.svg} \quad \texttt{M8.0\_positive\_bottom\_mirrored.svg}\\
\texttt{M8.0\_negative\_top.svg} \quad \texttt{M8.0\_negative\_bottom\_mirrored.svg}
\end{center}

\textbf{Polarity} (positive / negative) depends on your dry-film resist chemistry --- whether UV
exposure hardens the film (negative resist: you want opaque \emph{dots} where holes go, so the
resist there stays soft and washes out) or softens it (positive resist: the opposite). Check your
film's datasheet. Getting this wrong etches everything except your holes.

\textbf{Mirroring} is geometry, not chemistry. The two masks contact \emph{opposite faces} of the
same shim. Viewed through the shim they must be mirror images, or the two etch fronts will not
meet at the same $(x,y)$ and you get tapered, offset, or partially-blind holes. The bottom mask
is pre-mirrored in the file for exactly this reason.

\begin{warnbox}{Known gap in the generated files}
Polarity and side are \textbf{not distinguishable after the transparency is cut out}. Once you
have four unlabelled squares of film on the bench, you cannot tell which is which. Use
\texttt{combine\_masks\_for\_printing.py}, or mark each film in the margin immediately after
cutting and before it leaves the printer. A surgical fix to \texttt{generate\_masks.py} exists
but has not been applied.
\end{warnbox}

\section{Registration features already in every mask}

\begin{description}[leftmargin=0.8cm,style=nextline]
\item[Fiducial dots, $r = 0.75$\,mm, at the corners.] Align top and bottom masks by
superimposing these before clamping. If the fiducials do not coincide, the holes will not either.
\item[Calibration square, exactly 10.0\,mm.] Measure it with callipers on the \emph{printed}
transparency. If it is not 10.00\,mm, your printer has scaled the artwork and every hole diameter
is wrong by that same factor. This is the single most common way a photochemical etch fails, and
it fails silently.
\item[Blank margin, 8\,mm] around the array's bounding circle, with the blank size rounded up to a
clean 5\,mm increment. For $M=8$: $r_{\max} = 14.678$\,mm, giving a 50\,mm blank.
\end{description}

\begin{keybox}{Print discipline}
Print at 100\%, ``actual size,'' no fit-to-page, no margins adjustment. Verify with the 10\,mm
square \emph{every time}, including reprints --- driver settings do not always persist. A
laser-printed transparency at high toner density is standard; if the holes look grey rather than
solid black under a bright light, print again or double up two films in register.
\end{keybox}

\chapter{The test coupon: do this before the real shim}

The undercut compensation is a guess. One coupon converts it into a measurement, and costs a
scrap of the same shim stock.

\section{Coupon design}

On a small offcut of the \emph{same} 0.05\,mm brass, from the \emph{same} sheet, etch a row of
holes spanning the three compensation brackets plus two extra points:
\begin{center}
mask apertures of 300, 325, 350, 375, 400\,$\mu$m --- at least three holes each
\end{center}
Use the same resist, same etchant batch, same temperature, same agitation, same duration you
intend to use on the real part. Anything you vary later invalidates the calibration.

\section{Measurement}

Measure the resulting hole diameters. Options in rough order of preference:
\begin{itemize}[leftmargin=1.2cm]
\item USB microscope or phone macro with a stage micrometer / steel rule in the same frame
\item Backlit transmission photograph, hole diameter in pixels scaled by a known reference
\item Pin gauges, if you have them fine enough
\end{itemize}

For photographic measurement, apply the same discipline established in the earlier photography
rounds: \textbf{lock autofocus and auto-exposure manually}, shoot raw (DNG), convert to TIFF.
JPEG conversion --- including via \texttt{sips} --- bakes in a tone curve that shifts apparent
edge position and corrupts the measurement.

\section{Extracting the real undercut}
\begin{equation}
C_{\text{measured}} = \overline{d_{\text{etched}}} - d_{\text{mask}}
\end{equation}
Compute it at each of the five mask sizes. Two things to check:

\begin{enumerate}[leftmargin=1.2cm]
\item \textbf{Is $C$ constant across mask size?} It should be, to first order --- undercut is set by
etch depth and time, not by aperture width. If $C$ grows with aperture, your etch is running long
and small holes are being reached later; shorten the time or improve agitation.
\item \textbf{Where does $C$ sit in the $25$--$75\,\mu$m bracket?} Regenerate the masks with
$C_{\text{measured}}$ substituted for \texttt{UNDERCUT\_NOMINAL\_UM} before committing the real shim.
\end{enumerate}

\begin{keybox}{Why this is worth a slide, not just a step}
Reporting a measured undercut converts ``we used a published design rule'' into ``we calibrated
our process.'' It is cheap, it is genuinely quantitative, and it directly parallels the
calibration argument you are making about $c_{\min}$ --- showing you distinguish inherited
parameters from measured ones. That distinction is the thing IYPT judges reward.
\end{keybox}

\chapter{The etch itself}

\section{Process outline}

\begin{enumerate}[leftmargin=1.2cm]
\item \textbf{Clean.} Degrease the shim thoroughly --- isopropanol, then a fine abrasive scrub,
then rinse and dry completely. Any fingerprint becomes a resist adhesion failure and then a
random hole.
\item \textbf{Laminate.} Apply dry-film photoresist to both faces. Heat-laminate slowly; trapped
bubbles become pinholes in the resist, which become unwanted etched pits.
\item \textbf{Register and expose.} Align top and bottom masks using the corner fiducials, clamp
between glass, expose both faces. Exposure time comes from the film datasheet --- run a step
wedge if you have never used this film before.
\item \textbf{Develop.} Sodium carbonate solution, gentle agitation, until the pattern is crisp.
Under-development leaves resist scum in the holes and they will not etch through.
\item \textbf{Inspect before etching.} This is a free checkpoint. Hold to the light: every hole
should be clear resist-free brass, all $N$ of them, in the right lattice.
\item \textbf{Etch.} Ferric chloride, both faces, agitated. Warm etchant is faster but undercuts
proportionally more --- keep temperature consistent with the coupon run or the calibration
transfers badly.
\item \textbf{Stop and strip.} Rinse immediately and thoroughly when holes break through; every
extra second is extra undercut. Strip remaining resist.
\end{enumerate}

\begin{warnbox}{Safety --- not optional}
Ferric chloride stains permanently, is corrosive to skin and eyes, and attacks most metals
including sinks and drains. Nitrile gloves, eye protection, ventilation, plastic or glass vessels
only. Do not pour spent etchant down a drain: it is a dissolved-copper waste stream. Neutralise
and dispose according to local rules, or return it to a supplier that accepts it back.
\end{warnbox}

\section{Post-etch verification, in order}

\begin{description}[leftmargin=0.8cm,style=nextline]
\item[Hole count.] Count them. Compare against \texttt{N\_actual} in the coordinate CSV (50 for
$M=8$). A missing hole is a resist defect; a hole that should not exist is a resist pinhole.

\item[Diameter.] Measure a sample of holes, ideally including centre and edge of the array.
Compare against the 400\,$\mu$m target. Edge-versus-centre differences reveal uneven agitation.

\item[Pitch.] Measure the nearest-neighbour spacing directly. It must be 4000\,$\mu$m. If it is
not, your print was scaled --- go back to the 10\,mm calibration square.

\item[Blank size and $r_{\max}$.] For $M=8$: 50\,mm blank, $r_{\max} = 14.678$\,mm to the
outermost hole centre.

\item[Hole quality.] Inspect under magnification for taper (the two etch fronts missed each
other), ragged edges (over-etch), or partial breakthrough. A tapered hole is not the circular
aperture the model assumed and will not perform as predicted.
\end{description}

\chapter{What to measure with the finished array}

Fabrication is not the deliverable. A measured optical result is. Scope this tightly.

\section{Priority 1: acuity delta (highest value per hour)}

\begin{enumerate}[leftmargin=1.2cm]
\item Printed logMAR or Snellen chart at a fixed, \emph{measured} distance (2--3\,m).
\item Three conditions: bare uncorrected eye; single 400\,$\mu$m pinhole; the fabricated array.
\item Record the smallest resolvable line per condition. Three repeats, randomised order.
\item Fix and record $\dgap$ with a spacer. Record the subject's refractive error in dioptres,
and the ambient lighting.
\end{enumerate}

Report the \textbf{delta in logMAR}, not an absolute clinical claim. One honest sentence ---
``the array improved uncorrected acuity by $X$ logMAR at $M = \_\_$\,D'' --- is what turns this
from a simulation into an answer to the problem statement's actual verb, \emph{correct}.

\section{Priority 2: the array-versus-single-hole comparison}

This is the measurement that tests \emph{your} architecture specifically. If one hole illuminates
the pupil at a time, then \textbf{on-axis acuity through the array should match the single
pinhole}, while the array should provide usable vision across a much wider gaze range. Test both
halves:
\begin{itemize}[leftmargin=1.2cm]
\item On-axis: array $\approx$ single hole? (predicted yes)
\item Off-axis, eye rotated $10$--$20^\circ$: single hole goes dark, array does not. (predicted)
\item Any doubled or ghosted image at any gaze angle? (predicted no --- and if you \emph{do} see
one, that is a finding worth more than confirming the prediction)
\end{itemize}

\section{Priority 3: blur-versus-defocus, extending ExP3}

The earlier frosted-screen series gave clean monotonic blur growth (+39\% over 20\,mm screen
displacement) at 1000\,$\mu$m --- qualitatively consistent with the model but uncalibrated in
dioptres. Repeating it through the fabricated array, with the screen displacement converted to an
equivalent defocus in dioptres, upgrades that from qualitative to quantitative.

\begin{keybox}{Rig discipline, carried over from the earlier rounds}
Manual autofocus and auto-exposure lock on any camera --- the unlocked AF is the most likely
cause of the earlier blur-plateau anomaly. Shoot DNG, convert to TIFF, never JPEG for anything
you intend to measure.
\end{keybox}

\section{Report the negatives}

If the array underperforms the single hole, if the blind fraction is visible in use, if
transmission is uncomfortably dark at 1\% --- report it. This project's record already includes
acuity-model failures, a luminance re-sweep that made things worse, and two formula errors, all
documented rather than tuned away. That discipline is the most defensible thing you have. Do not
abandon it at the last step.

\appendix

\chapter{Quick-reference formula card}

\begin{longtable}{p{5.2cm}p{6.2cm}p{3.0cm}}
\caption{Every working formula in one place.}\\
\toprule
\textbf{Quantity} & \textbf{Formula} & \textbf{Status} \\
\midrule
\endhead
Myopic axial length & $L_{\text{ax}} = n'/(P_0-M)$ & live \\
Emmetropic axial length & $L_0 = n'/P_0 = 22.267$\,mm & live \\
ASM propagator & $H = \exp(ikz\sqrt{1-\lambda^2(f_x^2+f_y^2)})$ & live, used twice \\
Fresnel number & $N_F = R_{\text{hole}}^2/(\lambda z)$ & live (gate) \\
Lens phase & $\exp(-i\frac{k}{2f}(x^2+y^2))$ & live \\
Airy intensity & $I = I_0(2J_1(x)/x)^2$ & validation only \\
Encircled energy & $EE = 1-J_0^2(x)-J_1^2(x)$ & validation only \\
On-axis Fresnel & $I(0,z)=4I_0\sin^2(\pi N_F/2)$ & live (cross-check) \\
Relative illuminance & $E_{\text{rel}}=(\min(\ahole,\Dpupil)/\Dpupil)^2$ & live \\
Hole--pupil overlap test & $d-R_h<R_p$ & live \\
Overlap area & Eq.~\ref{eq:overlap} & live \\
Gaze pupil shift & $\Delta=R_{\text{rot}}\sin\varphi$, $R_{\text{rot}}\approx13$--14\,mm & live \\
Ghost separation & $\Delta r = sM/(P_0-M)$ & live, $R^2=0.995$ \\
Packing density area & $A_{\text{hull}}+P_{\text{hull}}R_h+\pi R_h^2$ & live \\
Undercut compensation & $d_{\text{mask}}=d_{\text{target}}-C$, $C\in[0.5t,1.5t]$ & live, uncalibrated \\
Geometric blur & $r_{\text{geo}}=\frac{a_{\text{stop}}}{2}\frac{|\vr-\vf|}{\vf}$ & \textbf{deprecated} \\
RSS combination & $\sqrt{r_{\text{geo}}^2+r_{\text{diff}}^2}$ & \textbf{deprecated} \\
Quadrature $a_{\text{eff}}$ & $\sqrt{\ahole^2+(2.44\lambda\dgap/\ahole)^2}$ & \textbf{deprecated} \\
\bottomrule
\end{longtable}

\chapter{Constants used throughout}

\begin{longtable}{p{4.5cm}p{3.5cm}p{6.0cm}}
\toprule
\textbf{Symbol} & \textbf{Value} & \textbf{Meaning} \\
\midrule
\endhead
$P_0$ & 60.0\,D & emmetropic eye power, held fixed \\
$n'$ & 1.336 & vitreous refractive index \\
$L_0 = f'$ & 22266.667\,$\mu$m & emmetropic axial length \\
$\Dpupil$ & 4000\,$\mu$m & default pupil diameter \\
$R_{\text{rot}}$ & 13--14\,mm & eye rotation radius \\
$\lambda$ & 550\,nm & default design wavelength \\
$\dgap$ & 12--15\,mm & spectacle vertex distance \\
$c_{\min}$ & 55\,$\mu$m & Shen fusion floor --- inherited, non-binding here \\
$t$ & 50\,$\mu$m & brass shim thickness \\
$C_{\text{undercut}}$ & 25 / 50 / 75\,$\mu$m & low / nominal / high compensation \\
Fiducial radius & 0.75\,mm & corner registration dots \\
Calibration square & 10.0\,mm & print scale verification \\
Blank margin & 8.0\,mm & around array bounding circle \\
\bottomrule
\end{longtable}

\chapter{Pre-fabrication checklist}

\begin{enumerate}[leftmargin=1.2cm]
\item Confirm the target design row in \texttt{gaze\_weighted\_winner\_per\_M.csv} and that
\texttt{fabrication\_summary.csv} matches it.
\item Open the coordinate CSV; confirm $N$, and hand-compute the nearest-neighbour pitch on two
holes as an independent check of the generator.
\item Confirm resist polarity against the film datasheet. Choose positive or negative masks.
\item Print all four masks; \textbf{measure the 10\,mm calibration square on each}.
\item Label each film in the margin before cutting --- polarity and side.
\item Run the test coupon. Measure $C_{\text{measured}}$.
\item Regenerate masks with $C_{\text{measured}}$ if it differs materially from 50\,$\mu$m.
\item Etch the real shim. Verify count, diameter, pitch, $r_{\max}$, hole quality.
\item Photograph the finished array against a scale, DNG, for the presentation.
\end{enumerate}

\end{document}
